\section{Related Work}\label{sec:related-work}
\textbf{Sycophancy in Psychology.}
Behaviors related to sycophancy have long been studied across several research areas in social psychology, including impression management \cite{goffman1955facework}, social conformity \cite{asch1956studies}, ingratiation \cite{jones1964ingratiation}, and social influence \cite{cialdini1984influence, cialdini2016presuasion}, each capturing related but distinct aspects of the phenomenon. 
Goffman \cite{goffman1955facework} characterized social interactions as pervasively organized by face-work and Jones \cite{jones1964ingratiation} formalized ingratiation as a reward-seeking class of tactics comprising of: other-enhancement, opinion conformity, and self-presentation. 
Cialdini~\cite{cialdini1984influence, cialdini2016presuasion} in turn distilled the broader compliance literature into seven heuristic principles --- reciprocity, commitment, social proof, liking, authority, scarcity, and unity ---each operating as a cue that elicits agreement under low cognitive engagement. 
Whereas Jones characterizes the strategic goals of the influencer, Cialdini characterizes the situational levers through which compliance is elicited from the target.
\cite{rehani2026} apply psychometric methodology from social psychology to the human-LLM setting and recovering three observer-distinguishable dimensions: uncritical agreement, obsequiousness, and excitement. 
Despite the rich literature on influence, most existing sycophancy benchmarks have drawn on it partially. 

\textbf{Sycophancy in LLMs.}
Sycophantic behavior was first documented systematically by \cite{perez2023modelwrittenevals}, who used model-written evaluations to show that LLMs trained with reinforcement learning human feedback  (RLHF) shift their opinions to align with user views. 
Various mitigation work have since been proposed, such as synthetic data augmentation~\cite{wei2023syntheticsycophancy}, direct preference optimization~\cite{khan2024mitigating}, pinpoint tuning~\cite{chen2024yes}, and probing \cite{papadatos2024probepenalty}, with~\cite{malmqvist2024sycophancymitigations} surveying their relative efficacy. 
On the evaluation side, \cite{sharma2023towards} introduced more realistic open-ended evaluation tasks with \textsc{SycophancyEval}. 
Subsequent sycophancy benchmarks generally fall into two categories. 
The first are domain specific, which includes mathematical theorem proving \cite{BROKENMATH}, medicine \cite{chen2025helpfulness}, and political opinion \cite{GermanPartiesQA}. 
The second are domain-general sycophancy benchmarks that isolate particular elicitation strategies \cite{SycEval, PARROT, BEACON, duffy2024sycobench}. 
However, across both groups, evaluations have been predominantly single turn, with recent benchmarks having started examining multi-turn dynamics \cite{TRUTHDECAY, FLIPFLOP, SYCONBENCH}. 
Beyond explicit change in opinion, \cite{ELEPHANT} introduce \textsc{ELEPHANT} which is grounded in social psychological theory to capture implicit forms of sycophantic behavior. 
A complementary line of work has shown that sycophancy is not fully captured by final-output scores: under social pressure, models exhibit post-hoc rationalization \cite{turpin2023languagemodels, turpin2023language, walden2026reasoning, chua2025reasoning}, with these influences typically not verbalized \cite{chua2025reasoning} and correct reasoning sometimes preceding sycophantic answers \cite{chang2026thermodynamic}.
This work brings together the multi-turn pressure dimension, grounded social dimensions of sycophancy, and the role of user influence as a trigger into a unified evaluation framework.

\textbf{Mechanistic Understanding of LLM Sycophancy.}
Several studies have shown that sycophancy is encoded in linear directions within the activation space \cite{jain2025psychometrictraits, papadatos2024probepenalty, chen2025personavectors, rimsky2024steering, vennemeyer2025causalsycophancy, genadi2026linear}. 
However, much of this literature has implicitly treated sycophancy as a single mechanism \cite{rimsky2024steering, chen2025personavectors, papadatos2024probepenalty}. 
Recent findings argue against this assumption, with sparse autoencoder analyses \cite{templeton2024scaling} revealing non-overlapping features associated with distinct sycophancy-related behaviors and steering interventions \cite{vennemeyer2025causalsycophancy} demonstrating that sycophantic agreement and praise occupy distinct subspaces that can be independently amplified or suppressed.
The conditions that trigger each subtype are less understood. \cite{wang2025truthoverridden} provide evidence that not all elicitation cues operate through the same mechanism: opinion
statements actively suppress the model's learned knowledge in later layers, while authority framing exerts negligible behavioral or representational effect. 
Chain-of-thought reasoning analysis shows that sycophantic commitment is not fixed at input but evolves dynamically across the reasoning trace~\cite{feng2026good, mirtaheri2026catching, duszenko2026anchors, hu2026MONICA}. 
Together, these findings motivate an evaluation framework that goes beyond measuring sycophancy as a single score, instead systematically varying triggers and examining how each modulates the reasoning process.

% \begin{table*}[t]
% \centering
% \begin{threeparttable}

% \caption{
%     Comparison of \ourWork and related datasets in physics. 
%     \textbf{Context}: Questions grounded in provided context or background information. \textbf{Context abbreviations}: (Hum = Humanities, SSci = Social Sciences, STEM, Med = Medicine).
%     \textbf{Trigger}: Questions that hinge on a specific prompt, cue, or perturbation. 
%     \textbf{Temporal}: Questions involving temporal or sequential patterns. \textbf{Trigger abbreviations}: I = Influence, A = Authority, SP = Social Proof,
%     C = Consistency, R = Reciprocity, L = Liking, S = Scarcity,
%     U = Unity, P = Plain Disagreement.
%     \textbf{Psych}: Questions grounded in established psychological or social-scientific theory.
% }
% \label{table:related_datasets}
% \begin{tabular}{@{}l|r|cccc
%     @{}} \toprule
%   & Size& Context & Trigger & Temporal & Psych\\ \midrule
% PARROT\cite{PARROT} & 11,040 & All & A & \cross & \cross \\
% Sycophancy Eval\cite{sharma2023towards} & 20,000+ & All & C/L/P/SP  & \cross  & \cross \\
% Persona Evals \cite{perez2023modelwrittenevals} & 30,000+ & SSci/STEM/Hum & L & \cross  & \cross \\
% SycEval\cite{SycEval} & 1,000 & STEM/Med & A/P & \cross  & \cross \\
% BrokenMath \cite{BROKENMATH} & 504 & STEM & C/SP & \cross  & \cross\\
% Scientific QA \cite{SCIENTIFICQA} & 1370? & STEM & P & \cross  & \cross \\
% Beacon \cite{BEACON} & 420 & SSci/Hum & P & \cross  & \cross \\
% DarkBench \cite{DARKBENCH} & 660 & STEM/SSci & L & \cross  & \cross \\
% Syco-Bench \cite{duffy2024sycobench} & 140 & All & P & \cross  & \cross\\
% GermanPartiesQA \cite{GermanPartiesQA} & 418 & SSci & L & \cross & \cross \\
% FLIPFLOP \cite{FLIPFLOP} & 1,352 & All & A/P & \cross   & \cross\\
% TRUTH DECAY \cite{TRUTHDECAY} & 800+ &  All & A/C/L/P/SP & \gcheck  & \cross \\
% SYCON Bench \cite{SYCONBENCH} & 2500 & All & A/C/L/P/SP & \gcheck &   \cross\\
% ELEPHANT\cite{ELEPHANT} & 10,404 & ? & ? & \cross  & \gcheck\\
% % Medical \cite{chen2025helpfulness} & & Med & A & \cross

% \midrule
% \ourWork  & \textbf{2,000} & \gcheck & \gcheck & \gcheck & \gcheck \\
% \bottomrule

% \end{tabular}
% \end{threeparttable}
% \vspace{-5pt}
% \end{table*}

\begin{table*}[t]
\centering
\begin{threeparttable}

\caption{
    Comparison of \ourWork and related sycophancy benchmarks. 
    \textbf{Psych}: Questions grounded in established psychological or social-scientific theory.
    \textbf{Context}: Questions grounded in provided context or background information. \textbf{Context abbreviations}: (Hum = Humanities, SSci = Social Sciences, STEM, Med = Medicine).
    \textbf{Trigger}: Questions that hinge on a specific prompt, cue, or perturbation. 
    \textbf{Temporal}: Questions involving temporal or sequential patterns. \textbf{Trigger abbreviations}: I = Influence, A = Authority, SP = Social Proof,
    C = Consistency, R = Reciprocity, L = Liking, S = Scarcity,
    U = Unity, P = Plain Disagreement.
    \textbf{GT}: Questions that contain a ground truth. 
}
\label{table:related_datasets}
\begin{tabular}{@{}l|r|cccc
    @{}} \toprule & 
  Psych & Context & Trigger & Temporal & GT\\ \midrule
PARROT\cite{PARROT} & \cross & All & A & \cross & GT \\
Sycophancy Eval\cite{sharma2023towards} & \cross & All & C/L/P/SP  & \cross & Both \\
Persona Evals \cite{perez2023modelwrittenevals} & \cross & SSci/STEM/Hum & L & \cross  & NGT \\
SycEval\cite{SycEval} & \cross & STEM/Med & A/P & \cross  & GT \\
BrokenMath \cite{BROKENMATH} & \cross & STEM & C/SP & \cross  & GT \\
Scientific QA \cite{SCIENTIFICQA} & \cross & STEM & P & \cross & GT \\
Beacon \cite{BEACON} & \cross & SSci/Hum & P & \cross  & NGT \\
DarkBench \cite{DARKBENCH} & \cross & STEM/SSci & L & \cross  & GT \\
Syco-Bench \cite{duffy2024sycobench} & \cross & All & P & \cross  & NGT \\
GermanPartiesQA \cite{GermanPartiesQA} & \cross & SSci & L & \cross & NGT  \\
FLIPFLOP \cite{FLIPFLOP} & \cross & All & A/P & \cross &  GT\\
TRUTH DECAY \cite{TRUTHDECAY} & \cross &  All & A/C/L/P/SP & \gcheck  & Both \\
SYCON Bench \cite{SYCONBENCH} & \cross & All & A/C/L/P/SP & \gcheck & Both  \\
ELEPHANT\cite{ELEPHANT} & \gcheck & ? & ? & \cross  & NGT \\
% Medical \cite{chen2025helpfulness} & & Med & A & \cross

\midrule
\ourWork  & \gcheck & All & All & \gcheck & Both \\
\bottomrule

\end{tabular}
\end{threeparttable}
\vspace{-5pt}
\end{table*}


